\documentclass[11pt,a4paper]{article}
\usepackage[margin=25mm,headheight=14pt]{geometry}
\usepackage[T1]{fontenc}
\usepackage[utf8]{inputenc}
\usepackage{lmodern,amsmath,amssymb}
\usepackage{xcolor,graphicx,booktabs,longtable,array}
\usepackage{microtype,enumitem,fancyhdr,titlesec}
\usepackage{hyperref}
\definecolor{accent}{HTML}{163B55}
\hypersetup{colorlinks=true,linkcolor=accent,urlcolor=accent,
 pdftitle={The square maximizes the first Steklov eigenvalue among parallelograms},pdfauthor={Matthew J. Colbrook},
 pdfsubject={AIM 022: mathematical result and proof}}
\urlstyle{same}
\Urlmuskip=0mu plus 1mu
\setlength{\parindent}{0pt}
\setlength{\parskip}{0.35em}
\setlength{\emergencystretch}{3em}
\setlist{itemsep=0.25em,topsep=0.3em}
\providecommand{\tightlist}{\setlength{\itemsep}{0pt}\setlength{\parskip}{0pt}}
\titleformat{\section}{\normalfont\large\bfseries\color{accent}}{\thesection}{0.7em}{}
\titleformat{\subsection}{\normalfont\normalsize\bfseries}{\thesubsection}{0.6em}{}
\pagestyle{fancy}
\fancyhf{}
\fancyhead[L]{\small\sffamily AIM 022}
\fancyhead[R]{\small\sffamily Matthew J. Colbrook}
\fancyfoot[C]{\thepage}
\renewcommand{\headrulewidth}{0.3pt}
\setcounter{secnumdepth}{0}
\allowdisplaybreaks[2]
\title{\textbf{The square maximizes the first Steklov eigenvalue among parallelograms}}
\author{Matthew J. Colbrook\\[3pt]
{\small Department of Applied Mathematics and Theoretical Physics}\\
{\small University of Cambridge, United Kingdom}\\
{\small\href{mailto:mjc249@maths.cam.ac.uk}{\texttt{mjc249@maths.cam.ac.uk}}}}
\date{3 November 2026}
\begin{document}
\maketitle
\begin{abstract}
We prove that among nondegenerate planar parallelograms with fixed perimeter, the square uniquely maximizes the first positive Steklov eigenvalue. The argument uses explicit trial functions and rigorous estimates for the constants in the square spectrum. The result addresses a restricted class in AIM 022; the polygonal optimization problem remains open.
\end{abstract}
\noindent\textbf{Repository target.}
This manuscript concerns AIM problem 022 at revision
\texttt{3fa50884d129}. Its scope is stated below.
\medskip

\textbf{Scope.} The following proves the perimeter-normalized Steklov
conjecture among all parallelograms, with strict rigidity. It does
\textbf{not} settle the full catalogued conjecture for arbitrary convex
quadrilaterals or for all polygons. A bounded literature search found
the known rectangle theorem and work on inverse Steklov spectra of
parallelograms, but I make no claim that this restricted result is new.

\hypertarget{022-theorem}{%
\section{Theorem}\label{022-theorem}}

Among all nondegenerate planar parallelograms of a given perimeter, the
square uniquely maximizes the first positive Steklov eigenvalue.

\hypertarget{022-notation-and-elementary-constants}{%
\subsection{Notation and elementary
constants}\label{022-notation-and-elementary-constants}}

Let \(Q=(-1,1)^2\), and let \(k\) be the unique root in \((0,\pi /2)\) of

\[
\tan k\,\tanh k=1.
\]

Put \(h=\tanh k\) and \(s=k h\). The first positive Steklov eigenvalue
of \(Q\) is \(s\); its eigenspace is generated by

\[
f(x,y)=\sin(kx)\cosh(ky),\qquad
g(x,y)=\cosh(kx)\sin(ky).
\]

The square's scale-invariant value is \(8s\). The square spectral
formula is standard; it also follows from the Robin separation argument
below.

Introduce

\[
\delta=\frac{k^2(1-h^4)+h^2}{2kh},\qquad
\beta=\frac hk-h^2,\qquad
\kappa=\frac{1+h^4}{2kh}.
\]

We need only the coarse inequalities

\[
\delta>\frac45,\qquad 0<\beta<\frac13,\qquad
\kappa>\frac23,\qquad
8s>\frac{224}{45}>2\sqrt3.
\tag{C}
\]

These are elementary rigorous inequalities, not numerical assumptions.
For example, Taylor series with their usual alternating or positive-tail
bounds give

\[
\frac{93}{100}<k<\frac{94}{100},\qquad
\frac{73}{100}<h<\frac{74}{100}.
\]

Substituting the appropriate endpoints in the displayed formulas proves
(C): use

\[
\delta\ge
\frac{(93/100)^2(1-(74/100)^4)+(73/100)^2}
{2(94/100)(74/100)},
\]

\[
\beta\le\frac{74}{93}-\left(\frac{73}{100}\right)^2,
\qquad
\kappa\ge\frac{1+(73/100)^4}{2(94/100)(74/100)},
\qquad
8s>8\frac{93}{100}\frac{73}{100}.
\]

The accompanying \texttt{parallelogram\_constant\_certificate.py} proves
every bracket and comparison by exact rational arithmetic. For
transparency, its exponential enclosure is the degree-20 Taylor sum plus
a geometric bound on the positive tail, while sine and cosine use
alternating sums through degree 17 and 16, respectively, followed by the
next negative term. The root function \(\sin x \tanh x-\cos x\) is
strictly increasing on \((0,\pi /2)\), so its rationally enclosed signs at
\(93/100\) and \(94/100\) certify the unique root bracket.

\hypertarget{022-parametrization}{%
\subsection{Parametrization}\label{022-parametrization}}

Scale the parallelogram so that its perimeter is eight. Translate its
center to the origin. After interchanging and reversing edge vectors
when necessary, it has the form

\[
P=\{xv+yw:(x,y)\in Q\},\qquad
|v|=a=1+t,\quad |w|=b=1-t,
\]

with \(0\le t<1\), and the angle between \(v,w\) belongs to \((0,\pi /2]\).
Set

\[
c=\cos\theta\in[0,1),\qquad z=\sin\theta=\sqrt{1-c^2}.
\]

The area is \(4abz\). Since the boundary is centrally symmetric, every
linear coordinate function centered at the origin has zero boundary
mean.

\hypertarget{022-two-linear-trial-estimates}{%
\subsection{Two linear trial
estimates}\label{022-two-linear-trial-estimates}}

For any vector \(e\), the boundary second moment is

\[
\int_{\partial P}(e\cdot X)^2\,ds
=\frac43(a+3b)(e\cdot v)^2
 +\frac43(3a+b)(e\cdot w)^2.
\tag{1}
\]

This follows by integrating separately along the four sides. Taking
\(e\) to be a unit vector parallel to the longer edge and dropping the
second nonnegative term gives

\[
8\sigma_1(P)
\le\frac{12(1-t)}{(1+t)(2-t)}.
\tag{2}
\]

Here the factor \(z\le 1\) was discarded. The right side decreases with
\(t\); for \(t\ge 1/8\), it is at most \(224/45<8s\). This proves strict
optimality outside \(t<1/8\).

Alternatively take \(e\) to be the unit angular bisector. Its scalar
products with \(v,w\) are \(a \cos (\theta /2), b \cos (\theta /2)\). Formula (1)
then yields

\[
8\sigma_1(P)
\le 6(1-t^2)\sqrt{\frac{1-c}{1+c}}
\le6\sqrt{\frac{1-c}{1+c}}.
\tag{3}
\]

Thus for \(c\ge 1/2\), the right side is at most \(2\sqrt3<8s\). It remains
to treat

\[
0\le t\le\frac18,\qquad 0\le c\le\frac12.
\tag{4}
\]

\hypertarget{022-exact-two-dimensional-trial-space}{%
\subsection{Exact two-dimensional trial
space}\label{022-exact-two-dimensional-trial-space}}

Pull the square functions \(f,g\) to \(P\) by the affine map
\((x,y)\mapsto xv+yw\). Both functions have boundary mean zero, and
their boundary cross term vanishes by parity.

For clarity define

\[
A_1=k^2\!\int_{-1}^1\!\cos^2(kx)dx
             \int_{-1}^1\!\cosh^2(ky)dy,
\qquad
A_2=k^2\!\int_{-1}^1\!\sin^2(kx)dx
             \int_{-1}^1\!\sinh^2(ky)dy,
\]

\[
d_a=2\cosh^2 k\int_{-1}^1\sin^2(kx)dx,
\qquad
d_b=2\sin^2 k\int_{-1}^1\cosh^2(ky)dy.
\]

Let \(E_0=A_1+A_2\) and \(B_0=d_a+d_b\). Direct integration, using
\(\tan k \tanh k=1\), gives

\[
E_0=sB_0,\quad
\frac{A_1-A_2}{E_0}=\delta,\quad
\frac{d_b-d_a}{B_0}=\beta,
\quad
\frac{\cosh(2k)-\cos(2k)}{E_0}=\kappa.
\]

Under the affine transformation, the energy and boundary-norm matrices
of the trial pair, normalized by \(E_0\) and \(B_0\), respectively, are

\[
\mathcal E=\frac1z
\begin{pmatrix}
\dfrac{1+t^2-2\delta t}{1-t^2}&-\kappa c\\[3pt]
-\kappa c&\dfrac{1+t^2+2\delta t}{1-t^2}
\end{pmatrix},
\qquad
\mathcal B=\begin{pmatrix}1-\beta t&0\\0&1+\beta t\end{pmatrix}.
\tag{5}
\]

The mixed energy uses

\[
\int_Q(f_xg_y+f_yg_x)\,dxdy
=\cosh(2k)-\cos(2k),
\]

whereas the other mixed integrals vanish by parity. Formula (5) is
exact.

\hypertarget{022-case-a-c8t2}{%
\subsection{\texorpdfstring{Case A:
\(c\ge 8t^2\)}{Case A: c >= 8t2}}\label{022-case-a-c8t2}}

Use the trial function \(f+g\). Formula (5) and (C) give

\[
\frac{\sigma_1(P)}s
\le\frac{(1+t^2)/(1-t^2)-\kappa c}{z}.
\tag{6}
\]

For \(t\le 1/8\),

\[
\frac{1+t^2}{1-t^2}\le1+\frac{128}{63}t^2.
\]

For \(c\le 1/2\),

\[
1-z\le(2-\sqrt3)c.
\]

Moreover

\[
\frac23-(2-\sqrt3)=\sqrt3-\frac43>\frac13.
\]

Therefore, if \(c\ge 8t^2\) and \((c,t)\ne (0,0)\),

\[
\kappa c-(1-z)>
\frac c3\ge\frac83t^2>\frac{128}{63}t^2
\]

when \(t>0\); for \(t=0,c>0\) the same strict comparison follows
directly. Hence the numerator in (6) is strictly smaller than \(z\), and
\(\sigma _1(P)<s\).

\hypertarget{022-case-b-c8t2}{%
\subsection{\texorpdfstring{Case B:
\(c<8t^2\)}{Case B: c\textless8t\^{}2}}\label{022-case-b-c8t2}}

This forces \(t>0\). Use \(f\) alone. Formula (5) gives

\[
\frac{\sigma_1(P)}s
\le\frac{1+t^2-2\delta t}
{z(1-t^2)(1-\beta t)}.
\tag{7}
\]

By (C), the factor excluding \(1/z\) is bounded above by

\[
\frac{1+t^2-\frac85t}{(1-t^2)(1-t/3)}<1-t.
\tag{8}
\]

To verify the strict inequality in (8), subtract its numerator from
\((1-t)(1-t^2)(1-t/3)\). The result is

\[
t\left(\frac4{15}-\frac53t+\frac43t^2-\frac13t^3\right)>0
\]

for \(0<t\le 1/8\), because the bracket is at least \(4/15-5/24-1/1536>0\).

Finally,

\[
z=\sqrt{1-c^2}\ge1-c^2>1-64t^4>1-t
\]

on this range. Inserting this into (7)--(8) proves \(\sigma _1(P)<s\).

The only remaining case is \(t=c=0\), which is the square and gives
equality. Scaling back proves the theorem.

\hypertarget{022-why-the-square-formula-is-the-first-eigenvalue}{%
\section{Why the square formula is the first
eigenvalue}\label{022-why-the-square-formula-is-the-first-eigenvalue}}

This paragraph can replace a citation if a self-contained account is
preferred. Consider the one-dimensional Robin operator on \((-1,1)\)
with outward boundary condition \(\partial_\nu u-su=0\). Since
\(0<s<1\), it has exactly one negative eigenvalue, \(-k^2\), with
eigenfunction \(\cosh (kx)\). Its first positive eigenvalue is \(k^2\),
with eigenfunction \(\sin (kx)\), because \(k \tanh k=k \cot k=s\). The
remaining eigenvalues exceed \(k^2\) by the one-dimensional Sturm
oscillation theorem.

On \(Q\), tensor products therefore give exactly one negative eigenvalue
of the Robin Laplacian, namely \(-2k^2\), and zero is its next
eigenvalue, of multiplicity two. Its zero eigenfunctions are \(f,g\)
above. Equivalently, the quadratic form
\(q(u)=\int_Q|\nabla u|^2-s\int_{\partial Q}u^2\) has negative index
one. If a mean-zero-boundary test function had \(q(u)<0\), then
\(q(1)<0\) and \(q(1,u)=-s\int_{\partial Q}u=0\), producing a
two-dimensional negative subspace, a contradiction. Thus all such
functions have Rayleigh quotient at least \(s\), while \(f,g\) attain
\(s\).

\hypertarget{022-sources-and-limitations}{%
\section{Sources and limitations}\label{022-sources-and-limitations}}

The benchmark square/rectangle formulas and general cuboid separation
method are treated in A. Girouard, J. Lagacé, I. Polterovich and A.
Savo, \emph{The Steklov spectrum of cuboids}, Mathematika \textbf{65}
(2019), 272--310, doi:10.1112/S0025579318000414; arXiv:1711.03075v2. The
present affine trial-space proof is given in full above, so its validity
does not depend on an unquoted theorem about general parallelograms.

The argument uses central symmetry twice: to ensure the trial functions'
boundary means vanish and to obtain the two-parameter affine
representation. A general quadrilateral has extra geometric parameters
and is not an affine square. Consequently this theorem is a restricted
result for catalogued problem 022, not a full resolution of that entry.
\end{document}
